\documentclass[12pt,a4paper]{article}

\usepackage[utf8]{inputenc}
\usepackage[T1]{fontenc}
\usepackage[norsk]{babel}
\usepackage{amsmath}
\usepackage{amssymb}
\usepackage{graphicx}
\usepackage{float}
\usepackage{booktabs}
\usepackage{hyperref}
\usepackage{cite}

\setlength{\parindent}{0pt}
\setlength{\parskip}{0.4em}

\title{LAB 3: Early Stage Assessment}
\author{Joakim Gual De Torrella Rivero Hopen}
\date{\today}

\begin{document}

\maketitle

% ============================================================
\section{Prosjekttittel}

\textbf{Analyse av økt kraftproduksjon i Nord-Norge og påvirkningen
på transmisjonsnettet mot Sør-Norge}


% ============================================================
\section{Problemstilling og mål}

\subsection{Problemstilling}

Prosjektet skal undersøke hvordan økt kraftproduksjon i Nord-Norge
påvirker kraftoverføringen gjennom transmisjonsnettet mot Sør-Norge.
Hovedfokuset er hvordan økt produksjon påvirker effektflyt,
belastning av overføringslinjer og spenningsforhold i nettet.
Det skal også undersøkes om begrensninger i overføringskapasiteten
kan redusere muligheten til å utnytte et eventuelt kraftoverskudd
i Nord-Norge.

Dette er relevant fordi et sterkere nett mellom Nord- og Sør-Norge
kan gjøre det mulig å utnytte kraftoverskuddet i nord bedre.
Bedre overføringskapasitet kan bidra til en mer effektiv utnyttelse
av kraftressursene og potensielt lavere strømpriser.

\subsection{Mål}

Målet med prosjektet er derfor å undersøke, gjennom simuleringer,
lastflytanalyse og data, hvordan økt kraftproduksjon i Nord-Norge
påvirker transmisjonsnettet mot Sør-Norge, og hvilke begrensninger
som eventuelt oppstår.

\subsection{Faglige spørsmål}

\begin{itemize}
    \item Hvordan påvirker økt kraftproduksjon i Nord-Norge
    effektflyten mot Sør-Norge?

    \item Kan økt kraftproduksjon føre til overbelastning
    av transmisjonslinjer?

    \item Hvordan påvirker økt kraftoverføring spenningene
    og effekttapene i nettet?

    \item Hvordan kan økt overføringskapasitet bidra til bedre
    utnyttelse av kraftproduksjonen i Nord-Norge?
\end{itemize}


% ============================================================
\section{Faglig grunnlag}

\subsection{Teori og faglig bakgrunn}

Det faglige grunnlaget for prosjektet vil hovedsakelig være
lastflytanalyse, overføringskapasitet, spenningsforhold og
effekttap i transmisjonsnettet. PowerFactory vil bli brukt
til å simulere hvordan økt kraftproduksjon påvirker nettet.

\subsection{Sentrale kilder}

\begin{itemize}
    \item \textbf{Fagbok:} \textit{Power System Analysis}
    av Grainger, Stevenson og Chang (2016).
    Til teori om lastflytanalyse, transmisjonslinjer
    og effekttap \cite{grainger2016}.

    \item \textbf{Vitenskapelig artikkel:}
    \textit{Spatio-temporal smoothing and dynamics of different
    electricity flexibility options for highly renewable energy
    systems---Case study for Norway} (2024).
    Til å undersøke hvordan kraftoverføring mellom regioner
    kan bidra til bedre utnyttelse av kraftproduksjonen
    \cite{grochowicz2024}.

    \item \textbf{Rapport:} Statnett,
    \textit{Analyse av transportkanaler 2025--2050}.
    Til informasjon om overføringskapasitet og begrensninger
    i det norske transmisjonsnettet \cite{statnett2025}.
\end{itemize}

\subsection{Litteratursøk}

Jeg planlegger å bruke ChatGPT til å finne relevant faglitteratur
og forskningsartikler. Deretter vil jeg kontrollere kildene og
vurdere hvor relevante de er for prosjektet. Jeg vil også bruke
Oria og Statnetts nettsider.


% ============================================================
\section{Foreløpig plan}

Jeg planlegger å undersøke tilgjengelige data om kraftproduksjon
og overføringskapasitet, og deretter bruke PowerFactory til å
simulere økt produksjon i Nord-Norge. Resultatene skal sammenlignes
for å undersøke hvordan effektflyten og belastningen i
transmisjonsnettet endres. Arbeidet vil bestå av litteratursøk,
modellering, simuleringer og analyse av resultatene fram mot
sluttrapporten.


% ============================================================
\section{Ambisjon og gjennomførbarhet}

Ambisjonen er å gjennomføre en realistisk analyse av hvordan
økt kraftproduksjon påvirker transmisjonsnettet. Den største
utfordringen blir trolig å finne tilstrekkelige data og lage
en representativ nettmodell. Det kan derfor bli nødvendig
å forenkle nettet og begrense analysen til utvalgte
overføringsforbindelser. Produksjonsnivå, linjeimpedanser
og overføringskapasitet forventes å ha stor innvirkning
på resultatene.


% ============================================================

\section{Bruk av kunstig intelligens}

ChatGPT fra OpenAI ble brukt som hjelpemiddel til litteratursøk,
kodegenerering og rettskrivning. Forslagene ble kontrollert
og tilpasset før bruk.

\section{GitHub-repository}

Prosjektfilene er tilgjengelige på GitHub:

\url{https://github.com/JoHopGDT/ELK330-LAB3}
\vspace{30mm}
% ============================================================
\bibliographystyle{IEEEtran}
\bibliography{references}

\end{document}